
\documentclass{article}
\usepackage{colortbl}
\usepackage{makecell}
\usepackage{multirow}
\usepackage{supertabular}

\begin{document}

\newcounter{utterance}

\centering \large Interaction Transcript for game `cladder', experiment `full\_v1.5\_default', episode 6304 with refineonly{-}Qwen3.5{-}9B.
\vspace{24pt}

{ \footnotesize  \setcounter{utterance}{1}
\setlength{\tabcolsep}{0pt}
\begin{supertabular}{c@{$\;$}|p{.15\linewidth}@{}p{.15\linewidth}p{.15\linewidth}p{.15\linewidth}p{.15\linewidth}p{.15\linewidth}}
    \# & \multicolumn{2}{c}{Player} && \multicolumn{2}{c}{Game Master} \\
    \hline

    \theutterance \stepcounter{utterance}  
    & & & \multicolumn{4}{p{0.6\linewidth}}{
        \cellcolor[rgb]{0.9,0.9,0.9}{
            \makecell[{{p{\linewidth}}}]{
                \texttt{\tiny{[P1$\langle$GM]}}
                \texttt{You are an expert in causal inference. The following question is not a typical commonsense query, but rather a meticulously designed question created by a professor specializing in causal inference, intended to assess the students' mastery of the course content.} \\
\texttt{Imagine a self{-}contained, hypothetical world with only the following conditions, and without any unmentioned factors or causal relationships: Zuph has a direct effect on rixq and xevu. Rixq has a direct effect on xevu. The overall probability of rixq is 32\%. The probability of not rixq and xevu is 35\%. The probability of rixq and xevu is 21\%. Is the chance of xevu smaller when observing rixq?} \\
\texttt{Start your answer with 'Yes' or 'No', followed by additional reasoning or evidence to support your explanation.} \\
            }
        }
    }
     \\ \\

    \theutterance \stepcounter{utterance}  
    & & \multicolumn{4}{p{0.6\linewidth}}{
        \cellcolor[rgb]{0.95,0.95,0.95}{
            \makecell[{{p{\linewidth}}}]{
                \texttt{\tiny{[P1$\leftrightarrow$P]}}
                \texttt{{[}DRAFT{]}} \\
\texttt{No} \\
\\ 
\texttt{**Reasoning:**} \\
\\ 
\texttt{To determine whether the chance of \$xevu\$ is smaller when observing \$rixq\$, we must compare the conditional probability \$P(xevu \textbackslash{}mid rixq)\$ with the marginal probability \$P(xevu)\$.} \\
\\ 
\texttt{1.  **Calculate the marginal probability \$P(xevu)\$:**} \\
\texttt{    We are given the probability of the joint event where \$rixq\$ is false and \$xevu\$ is true (\$P(\textbackslash{}neg rixq \textbackslash{}cap xevu) = 35\textbackslash{}\%\$) and the probability where both are true (\$P(rixq \textbackslash{}cap xevu) = 21\textbackslash{}\%\$).} \\
\texttt{    Since these two events (\$\textbackslash{}neg rixq \textbackslash{}cap xevu\$ and \$rixq \textbackslash{}cap xevu\$) are mutually exclusive and cover all cases where \$xevu\$ occurs:} \\
\texttt{    \$\$P(xevu) = P(\textbackslash{}neg rixq \textbackslash{}cap xevu) + P(rixq \textbackslash{}cap xevu)\$\$} \\
\texttt{    \$\$P(xevu) = 35\textbackslash{}\% + 21\textbackslash{}\% = 56\textbackslash{}\%\$\$} \\
\\ 
\texttt{2.  **Calculate the conditional probability \$P(xevu \textbackslash{}mid rixq)\$:**} \\
\texttt{    By definition, \$P(xevu \textbackslash{}mid rixq) = \textbackslash{}frac\{P(rixq \textbackslash{}cap xevu)\}\{P(rixq)\}\$.} \\
\texttt{    We are given \$P(rixq) = 32\textbackslash{}\%\$ and \$P(rixq \textbackslash{}cap xevu) = 21\textbackslash{}\%\$.} \\
\texttt{    \$\$P(xevu \textbackslash{}mid rixq) = \textbackslash{}frac\{21\textbackslash{}\%\}\{32\textbackslash{}\%\} = 0.65625 \textbackslash{}text\{ (or \} 65.625\textbackslash{}\%)\$\$} \\
\\ 
\texttt{3.  **Compare the probabilities:**} \\
\texttt{    *   \$P(xevu) = 56\textbackslash{}\%\$} \\
\texttt{    *   \$P(xevu \textbackslash{}mid rixq) \textbackslash{}approx 65.6\textbackslash{}\%\$} \\
\\ 
\texttt{The question asks if the chance of \$xevu\$ is **smaller** when observing \$rixq\$. Since \$65.6\textbackslash{}\% > 56\textbackslash{}\%\$, the chance of \$xevu\$ is actually **larger** when observing \$rixq\$. This is consistent with the causal structure described (\$Zuph \textbackslash{}to rixq \textbackslash{}to xevu\$), where \$rixq\$ acts as a positive cause for \$xevu\$.} \\
\\ 
\texttt{Therefore, the answer is **No**.} \\
            }
        }
    }
    & \\ \\

    \theutterance \stepcounter{utterance}  
    & \multicolumn{4}{p{0.6\linewidth}}{
        \cellcolor[rgb]{0.561,0.737,0.561}{
            \makecell[{{p{\linewidth}}}]{
                \texttt{\tiny{[P1$\rangle$GM]}}
                \texttt{No} \\
\\ 
\texttt{**Reasoning:**} \\
\\ 
\texttt{To determine whether the chance of \$xevu\$ is smaller when observing \$rixq\$, we must compare the conditional probability \$P(xevu \textbackslash{}mid rixq)\$ with the marginal probability \$P(xevu)\$.} \\
\\ 
\texttt{1.  **Calculate the marginal probability \$P(xevu)\$:**} \\
\texttt{    We are given the probability of the joint event where \$rixq\$ is false and \$xevu\$ is true (\$P(\textbackslash{}neg rixq \textbackslash{}cap xevu) = 35\textbackslash{}\%\$) and the probability where both are true (\$P(rixq \textbackslash{}cap xevu) = 21\textbackslash{}\%\$).} \\
\texttt{    Since these two events (\$\textbackslash{}neg rixq \textbackslash{}cap xevu\$ and \$rixq \textbackslash{}cap xevu\$) are mutually exclusive and cover all cases where \$xevu\$ occurs:} \\
\texttt{    \$\$P(xevu) = P(\textbackslash{}neg rixq \textbackslash{}cap xevu) + P(rixq \textbackslash{}cap xevu)\$\$} \\
\texttt{    \$\$P(xevu) = 35\textbackslash{}\% + 21\textbackslash{}\% = 56\textbackslash{}\%\$\$} \\
\\ 
\texttt{2.  **Calculate the conditional probability \$P(xevu \textbackslash{}mid rixq)\$:**} \\
\texttt{    By definition, \$P(xevu \textbackslash{}mid rixq) = \textbackslash{}frac\{P(rixq \textbackslash{}cap xevu)\}\{P(rixq)\}\$.} \\
\texttt{    We are given \$P(rixq) = 32\textbackslash{}\%\$ and \$P(rixq \textbackslash{}cap xevu) = 21\textbackslash{}\%\$.} \\
\texttt{    \$\$P(xevu \textbackslash{}mid rixq) = \textbackslash{}frac\{21\textbackslash{}\%\}\{32\textbackslash{}\%\} = 0.65625 \textbackslash{}text\{ (or \} 65.625\textbackslash{}\%)\$\$} \\
\\ 
\texttt{3.  **Compare the probabilities:**} \\
\texttt{    *   \$P(xevu) = 56\textbackslash{}\%\$} \\
\texttt{    *   \$P(xevu \textbackslash{}mid rixq) \textbackslash{}approx 65.6\textbackslash{}\%\$} \\
\\ 
\texttt{The question asks if the chance of \$xevu\$ is **smaller** when observing \$rixq\$. Since \$65.6\textbackslash{}\% > 56\textbackslash{}\%\$, the chance of \$xevu\$ is actually **larger** when observing \$rixq\$. This is consistent with the causal structure described (\$Zuph \textbackslash{}to rixq \textbackslash{}to xevu\$), where \$rixq\$ acts as a positive cause for \$xevu\$.} \\
\\ 
\texttt{Therefore, the answer is **No**.} \\
            }
        }
    }
    & & \\ \\

    \theutterance \stepcounter{utterance}  
    & & & \multicolumn{2}{p{0.3\linewidth}}{
        \cellcolor[rgb]{0.9,0.9,0.9}{
            \makecell[{{p{\linewidth}}}]{
                \texttt{\tiny{[GM$|$GM]}}
                \texttt{No} \\
            }
        }
    }
    & & \\ \\

    \theutterance \stepcounter{utterance}  
    & & & \multicolumn{2}{p{0.3\linewidth}}{
        \cellcolor[rgb]{0.9,0.9,0.9}{
            \makecell[{{p{\linewidth}}}]{
                \texttt{\tiny{[GM$|$GM]}}
                \texttt{no} \\
            }
        }
    }
    & & \\ \\

    \theutterance \stepcounter{utterance}  
    & & & \multicolumn{2}{p{0.3\linewidth}}{
        \cellcolor[rgb]{0.9,0.9,0.9}{
            \makecell[{{p{\linewidth}}}]{
                \texttt{\tiny{[GM$|$GM]}}
                \texttt{game\_result = WIN} \\
            }
        }
    }
    & & \\ \\

\end{supertabular}
}

\end{document}
